\chapter{Finestra Principale (ConfiguratorApp)}

\section{Scopo}
La finestra principale e il punto di ingresso del configuratore.
Da qui l'utente apre tutte le finestre di configurazione di dominio, meteo, dispersione e operazioni farm.

\section{Azioni disponibili}
\begin{itemize}[leftmargin=1.5em]
  \item Definisci Dominio
  \item Orografia e Uso Terreno
  \item Dominio Temporale
  \item Configura CALMET
  \item Configura CALPUFF
  \item Configurazione Farm
  \item Operazioni sul Farm
  \item Salva Configurazione
  \item Carica Configurazione
  \item Clear Simulazion
\end{itemize}

\section{Sottofinestre e dialog dal main}
\subsection{Dialog Salva Configurazione}
Il comando \emph{Salva Configurazione} richiede un nome tramite dialog di input.
Alla conferma copia i file \texttt{*.json} da \texttt{temp\_config/} in una nuova cartella timestampata in \texttt{saved\_configurations/} e crea \texttt{metadata.json}.

\subsection{Finestra Carica Configurazione}
Il comando \emph{Carica Configurazione} apre una finestra di selezione delle configurazioni salvate.
Il dettaglio funzionale e riportato nel Capitolo~\ref{chap:carica-config}.

\subsection{Dialog Conferma Pulizia}
Il comando \emph{Clear Simulazion} mostra una conferma e, in caso positivo, svuota i contenuti di:
\begin{itemize}[leftmargin=1.5em]
  \item \texttt{temp\_config/}
  \item \texttt{Outputs/}
  \item tutte le cartelle \texttt{*\_INP/}
\end{itemize}
Le cartelle vengono mantenute, ma il contenuto viene rimosso.

\section{Dipendenze file principali}
\begin{itemize}[leftmargin=1.5em]
  \item Lettura/scrittura configurazioni temporanee in \texttt{temp\_config/}
  \item Snapshot versionati in \texttt{saved\_configurations/}
  \item Gestione artefatti intermedi in \texttt{Outputs/} e \texttt{*\_INP/}
\end{itemize}
